\chapter{Introduction to GUI Programming in Python (Tkinter)}

\section{Introduction: Real-Life Hook}

Imagine this:  

\begin{itemize}
    \item You open a calculator app on your phone  
    \item Click buttons to enter numbers and see results instantly  
    \item Games, quizzes, and chat apps all have interactive windows  
\end{itemize}

In Python, \textbf{Tkinter} allows you to create \textbf{Graphical User Interfaces (GUIs)} --- windows, buttons, text boxes --- so programs become \textbf{visual and interactive} instead of just text-based.

\section{Concept: What is Tkinter?}

\begin{definition}{Tkinter}
Tkinter is Python’s standard library for \textbf{creating GUI applications}. It provides widgets like windows, labels, buttons, and entry fields.
\end{definition}

\section{Creating Your First Window}

\begin{verbatim}
import tkinter as tk

# Create main window
window = tk.Tk()
window.title("My First GUI")
window.geometry("300x200")  # width x height

# Run the GUI loop
window.mainloop()
\end{verbatim}

- \texttt{Tk()} creates the main window  
- \texttt{title()} sets the window title  
- \texttt{geometry()} sets window size  
- \texttt{mainloop()} keeps the window running  

\section{Adding Labels and Buttons}

\begin{verbatim}
import tkinter as tk

window = tk.Tk()
window.title("Simple GUI")
window.geometry("300x200")

# Add a label
label = tk.Label(window, text="Hello, Tkinter!", font=("Arial", 16))
label.pack(pady=20)

# Add a button
def greet():
    label.config(text="Welcome to GUI Programming!")

button = tk.Button(window, text="Click Me", command=greet)
button.pack(pady=10)

window.mainloop()
\end{verbatim}

- \texttt{Label} displays text  
- \texttt{Button} triggers a function when clicked  
- \texttt{command} links a function to the button  

\section{Mini-Project: Simple Calculator GUI}

Create a calculator that adds two numbers:

\begin{verbatim}
import tkinter as tk

def add_numbers():
    try:
        num1 = float(entry1.get())
        num2 = float(entry2.get())
        result_label.config(text=f"Result: {num1 + num2}")
    except ValueError:
        result_label.config(text="Enter valid numbers!")

window = tk.Tk()
window.title("Calculator")
window.geometry("300x200")

tk.Label(window, text="Number 1:").pack()
entry1 = tk.Entry(window)
entry1.pack()

tk.Label(window, text="Number 2:").pack()
entry2 = tk.Entry(window)
entry2.pack()

tk.Button(window, text="Add", command=add_numbers).pack(pady=10)

result_label = tk.Label(window, text="Result:")
result_label.pack()

window.mainloop()
\end{verbatim}

- Students learn \textbf{widgets, layout, and event handling}  

\section{Step-by-Step Diagram: GUI Event Flow}

\begin{center}
\begin{tikzpicture}[node distance=1.5cm]
\node (start) [draw, ellipse, fill=green!20] {Start GUI Program};
\node (create) [draw, rectangle, fill=yellow!20, below of=start] {Create Window and Widgets};
\node (event) [draw, diamond, aspect=2, fill=orange!20, below of=create, yshift=-0.5cm] {User Event Occurs (Click/Input)};
\node (action) [draw, rectangle, fill=yellow!20, below of=event, yshift=-0.5cm] {Execute Associated Function};
\node (update) [draw, rectangle, fill=orange!20, below of=action] {Update Widgets / Display Results};
\node (end) [draw, ellipse, fill=green!20, below of=update] {Program Ends};
\draw[->] (start) -- (create);
\draw[->] (create) -- (event);
\draw[->] (event) -- (action);
\draw[->] (action) -- (update);
\draw[->] (update) -- (end);
\end{tikzpicture}
\end{center}

\section{Key Concepts}

\begin{itemize}
    \item Tkinter is used to create \textbf{GUI applications}  
    \item \textbf{Widgets}: Label, Button, Entry, etc.  
    \item \textbf{Event-driven programming}: functions respond to user actions  
    \item \texttt{mainloop()} keeps the GUI running  
    \item GUI programming makes programs \textbf{interactive and user-friendly}  
\end{itemize}

\section{End-of-Chapter Exercises}

\begin{exercise}{Very Short Questions}
\begin{enumerate}
    \item What is Tkinter in Python?  
    \item Name three common Tkinter widgets.  
    \item What does \texttt{mainloop()} do?  
\end{enumerate}
\end{exercise}

\begin{exercise}{Short Answer Questions}
\begin{enumerate}
    \item Write a program to create a window with a label “Welcome!”  
    \item How do you link a button to a function in Tkinter?  
    \item Explain the difference between \texttt{pack()}, \texttt{grid()}, and \texttt{place()} layout managers.  
\end{enumerate}
\end{exercise}

\begin{exercise}{Activity-Based Questions}
\begin{enumerate}
    \item Create a GUI that takes two numbers and subtracts them when a button is clicked.  
    \item Create a quiz GUI with a question label, entry box for answer, and a button to check the answer.  
    \item Design a simple “To-Do List” GUI that allows users to add tasks and display them in a list.  
\end{enumerate}
\end{exercise}

\section{Teacher Notes}

\marginnote{
\textbf{Learning Objectives:}  
\begin{itemize}
    \item Students can create windows, labels, buttons, and entry fields  
    \item Students understand event-driven programming  
    \item Students can build simple interactive GUI applications  
\end{itemize}

\textbf{Teaching Tips:}  
\begin{itemize}
    \item Start with simple widgets and add functionality gradually  
    \item Encourage students to experiment with colors, fonts, and layout  
    \item Relate GUI elements to real-life applications (calculators, games)  
\end{itemize}

\textbf{Assessment:}  
\begin{itemize}
    \item Check GUI creation and widget functionality  
    \item Observe event handling and interactive response  
\end{itemize}
}